\chapter{봄베: \texorpdfstring{\code{bombe}}{bombe}}\label{ch:rust-bombe}

이 장이 2부의 중심이다. \code{bombe.rs}는 봄베 한 대를 옮긴 것이고, 3부의 회로는 이 파일과 같은 답을 내야
한다. 읽는 순서는 상태, 스크램블러 표, 전류의 확산, 구동이다.

\section{상태: 정수 26개}

\ref{sec:bombe-one-page}절에서 본 대로 봄베의 상태는 26개 케이블 × 26가닥 = 676가닥의 통전 여부다.
코드에서는 이것이 다음 한 줄이다.

\begin{lstlisting}[language=Rust,numbers=none]
let mut live = [0u32; N];
\end{lstlisting}

32비트 정수 26개의 배열이다. \code{live[c]}가 케이블 $c$이고, 그 정수의 $w$번째 비트가 1이면 케이블 $c$의
$w$번 선이 통전된 것, 곧 ``글자 $c$의 스테커 짝은 $w$다''라는 가설이 살아 있는 것이다
(그림~\ref{fig:g-state}).

\begin{figure}[ht]
\centering
\begin{tikzpicture}[x=3.6mm,y=-4.4mm]
  \foreach \r/\lab/\bits in {0/A/{16},1/B/{},2/C/{},3/D/{},4/E/{0,1,2,3,4,5,6,7,8,9,10,11,12,13,14,15,16,17,18,19,20,22,23,24,25}} {
    \node[font=\ttfamily\footnotesize,anchor=east] at (-0.8,\r) {live[\r]};
    \node[font=\footnotesize,anchor=east,text=cgray] at (-5.2,\r) {케이블 \texttt{\lab}};
    \foreach \c in {0,...,25} { \draw[cgray!60,fill=white] (\c-0.5,\r-0.42) rectangle (\c+0.5,\r+0.42); }
    \foreach \c in \bits { \fill[cblue] (\c-0.5,\r-0.42) rectangle (\c+0.5,\r+0.42); }
  }
  \node at (12.5,5.1) {$\vdots$};
  \foreach \l [count=\i from 0] in {a,b,c,d,e,f,g,h,i,j,k,l,m,n,o,p,q,r,s,t,u,v,w,x,y,z} {
    \node[font=\ttfamily\tiny] at (\i,-0.95) {\l};
  }
  \node[note,anchor=west] at (26.3,0) {$\Stk(\mathtt{A})=\mathtt{Q}$ 하나};
  \node[note,anchor=west,text=corange] at (26.3,4) {테스트 레지스터: 25/26};
  \draw[corange,line width=0.9pt] (-0.5,3.58) rectangle (25.5,4.42);
\end{tikzpicture}
\caption{\code{live} 배열(일부). 가로줄 하나가 정수 하나, 칸 하나가 비트 하나이자 전선 한 가닥이다.
테스트 글자가 \texttt{E}이면 \code{live[4]}가 테스트 레지스터다.}
\label{fig:g-state}
\end{figure}

\Why 26가닥을 정수 하나에 담으면 두 가지가 싸진다. 선 하나를 통전시키는 것은 비트 하나를 켜는 것
(\code{live[c] |= 1 << w})이고, ``26가닥이 전부 통전되었는가''는 정수 비교 한 번(\code{live[c] == FULL})이다.

\General 비트 하나에 참/거짓 하나를 담는 이 방법은 흔하다. 파일의 권한(읽기, 쓰기, 실행)을 수 하나의 비트로
나타내는 것, 요일 일곱 개의 선택을 7비트로 나타내는 것이 같은 생각이다.

\Ours 이 표현은 회로와 일대일로 대응한다. 3부의 \code{bombe\_core}에는 676비트 레지스터 \code{live\_q}가
있고, 그 $26c + w$번째 비트가 여기의 \code{live[c]}의 $w$번째 비트다. 검증할 때에는 이 676비트를 통째로
비교한다.

\section{스크램블러 표와 위치}

\begin{lstlisting}[language=Rust,caption={\code{bombe.rs}: 드럼 위치와 스크램블러 표(주석 생략, 그리스 로터 부분 생략).},label={lst:scramblers}]
pub struct Position {
    pub greek: u8,
    pub left: u8,
    pub middle: u8,
    pub right: u8,
}

fn inner(&self, g: u8, l: u8, m: u8) -> [u8; N] {
    let [lr, mr, _] = self.rotors;
    std::array::from_fn(|c| {
        let mut x = through(&mr.forward, m, c as u8);
        x = through(&lr.forward, l, x);
        ...
        x = self.reflector.map[x as usize];
        ...
        x = through(&lr.backward, l, x);
        through(&mr.backward, m, x)
    })
}

pub fn scramblers(&self, g: u8, l: u8, m: u8) -> [[u8; N]; N] {
    let inner = self.inner(g, l, m);
    let rr = self.rotors[2];
    std::array::from_fn(|r| {
        std::array::from_fn(|c| {
            let x = through(&rr.forward, r as u8, c as u8);
            through(&rr.backward, r as u8, inner[x as usize])
        })
    })
}
\end{lstlisting}

\code{Position}은 드럼의 위치, 정확히는 \emph{크립의 첫 글자를 암호화하는 순간}의 네 로터의 코어 위치다
(3로터 기종에서는 \code{greek}이 0). \code{std::array::from\_fn(|c| ...)}은 ``번호 \code{c}마다 이 식을 계산해
배열을 만든다''는 뜻이다.

\code{inner}는 스크램블러의 안쪽, 곧 가운데 로터에서 반사판을 거쳐 다시 가운데 로터까지의 치환을 26글자에
대해 계산한다. \code{scramblers}는 그 안쪽 치환의 앞뒤에 오른쪽 로터를 씌우되, 오른쪽 로터의 26가지 위치
\code{r}에 대해 모두 계산해 표 26개를 만든다.

\Why 표를 이렇게 나누는 이유는 무엇이 자주 바뀌고 무엇이 드물게 바뀌는가에 있다. 메뉴의 모든 스크램블러는
가운데와 왼쪽 드럼의 위치가 같고, 오른쪽 드럼만 크립 위치 $k$만큼 어긋나 있다. 크립 위치 $k$의 선에 필요한
것은 오른쪽 코어 위치 $r+k$의 스크램블러다. 그러므로 가운데와 왼쪽이 정해진 동안 필요한 표는 이 26개가
전부이고, 선은 그 가운데 $(r+k) \bmod 26$번째를 골라 쓰기만 하면 된다. 오른쪽 드럼이 26칸을 도는 동안
표를 다시 계산할 일이 없다.

\General 구구단을 외워 두고 곱셈을 할 때마다 표에서 꺼내 쓰는 것과 같다. 같은 계산을 여러 번 해야 하면
한 번 해 두고 꺼내 쓴다. 다만 표가 쓸모 있으려면 표를 만드는 데 쓴 조건(여기서는 가운데와 왼쪽의 위치)이
한동안 변하지 않아야 한다.

\Ours \ref{sec:tradeoff}절에서 본 FPGA 엔진의 구조가 이것이다. 엔진은 가운데 로터가 바뀔 때만 이 표
26개를 계산해 메모리에 넣는다. 소프트웨어에서 속도를 위해 한 일이 회로에서는 크기를 줄이는 방법이 된다.

\section{메뉴를 기계에 걸기: \texorpdfstring{\code{Plugged}}{Plugged}}

\begin{lstlisting}[language=Rust,caption={\code{bombe.rs}: 배선된 메뉴.}]
pub struct Plugged {
    /// For each cable: (other cable, edge index).
    adj: Vec<Vec<(u8, usize)>>,
    /// Crib position of each edge.
    edge_pos: Vec<u16>,
    pub test_letter: u8,
    pub input_wire: u8,
}
\end{lstlisting}

\code{Plugged}는 메뉴를 봄베에 배선한 상태다. 실제 조작자가 뒷면의 잭에 케이블을 꽂던 일에 해당한다.
\code{adj[a]}는 글자 $a$의 케이블에 이어진 스크램블러들의 목록이다. 항목 하나는 ``건너편 글자, 선의 번호''다.
메뉴의 선 하나는 양쪽 글자의 목록에 한 번씩 들어간다. 전류가 양방향으로 흐르기 때문이다.

\section{전류의 확산: \texorpdfstring{\code{propagate}}{propagate}}\label{sec:propagate}

이 함수가 봄베다.

\begin{lstlisting}[language=Rust,caption={\code{bombe.rs}: 전류의 확산.},label={lst:propagate}]
pub fn propagate(&self, tables: &[[u8; N]], edge_table: &[usize],
                 start: (u8, u8), diagonal: bool, stop_when_full: bool) -> [u32; N] {
    let mut live = [0u32; N];
    let mut stack: [(u8, u8); N * N] = [(0, 0); N * N];
    let mut sp = 0usize;
    let test = self.test_letter as usize;
    live[start.0 as usize] |= 1 << start.1;
    stack[sp] = start;
    sp += 1;
    while sp > 0 {
        sp -= 1;
        let (a, w) = stack[sp];
        for &(b, e) in &self.adj[a as usize] {
            let w2 = tables[edge_table[e]][w as usize];
            let bit = 1u32 << w2;
            if live[b as usize] & bit == 0 {
                live[b as usize] |= bit;
                stack[sp] = (b, w2);
                sp += 1;
                if stop_when_full && b as usize == test && live[test] == FULL {
                    return live;
                }
            }
        }
        if diagonal {
            let bit = 1u32 << a;
            if live[w as usize] & bit == 0 {
                live[w as usize] |= bit;
                stack[sp] = (w, a);
                sp += 1;
                if stop_when_full && w as usize == test && live[test] == FULL {
                    return live;
                }
            }
        }
    }
    live
}
\end{lstlisting}

한 줄씩 읽는다.

\begin{itemize}
\item 3--9행. 상태를 모두 0으로 하고, 전압을 거는 선 \code{start}(케이블, 선) 하나를 켠다. \code{stack}은
``새로 통전되었지만 아직 이웃을 살피지 않은 선''의 목록이다. 선은 676가닥뿐이고 한 선은 한 번만 들어가므로
크기 676이면 넘치지 않는다.
\item 10--12행. 목록이 빌 때까지, 선 하나 $(a, w)$를 꺼낸다. ``글자 $a$의 짝이 $w$''라는 가설이다.
\item 13--24행. \textbf{스크램블러.} 케이블 $a$에 이어진 스크램블러마다, 건너편 케이블 $b$의 어느 선으로
가는지 표에서 읽는다(\code{w2}). ``$a$의 짝이 $w$이면 $b$의 짝은 \code{w2}''라는 추론이다. 그 선이 아직
꺼져 있으면 켜고 목록에 넣는다.
\item 25--35행. \textbf{대각선 보드.} 케이블 $a$의 $w$번 선이 켜졌으면 케이블 $w$의 $a$번 선을 켠다.
``$a$의 짝이 $w$이면 $w$의 짝은 $a$''다. 이 여섯 줄이 웰치먼의 개량 전부다.
\item 20--22행과 31--33행. \textbf{조기 종료.} 테스트 레지스터의 26가닥이 다 켜지면 그 위치는 정지가 아닌 것이
확정되었으므로 곧바로 돌아간다.
\end{itemize}

\Why 이 함수가 구하는 것은 ``전압을 건 선에서 배선을 따라 닿을 수 있는 선 전체''다. 새로 닿은 선을 목록에
넣고, 하나씩 꺼내 그 이웃을 살피는 방식은 이 집합을 빠짐없이 찾는다. 닿을 수 있는 선은 언젠가 목록에
들어가고, 목록에 들어간 선은 언젠가 꺼내져 이웃이 살펴지기 때문이다. 한 번 켜진 선은 다시 꺼지지 않고
(\code{if ... == 0}으로 이미 켜진 선은 건너뛴다) 선은 676가닥뿐이므로 반드시 끝난다. 그리고 결과는 선을
꺼내는 순서와 무관하다. 닿을 수 있는 선의 집합은 하나뿐이기 때문이다.

\General 그림판의 ``채우기'' 기능이 같은 알고리즘이다. 한 점에서 시작해 같은 색으로 이어진 점을 모두 칠한다.
새로 칠한 점을 목록에 넣고, 하나씩 꺼내 그 상하좌우를 살핀다. 소문이 퍼지는 것도 같다. 소문을 들은 사람이
아는 사람 모두에게 전하고, 이미 들은 사람은 다시 전하지 않는다.

\Ours 실제 봄베는 이 계산을 전류로 한다. 전선은 이어져 있기만 하면 전압이 저절로 전해지므로 ``목록''도
``순서''도 없다. 회로(3부)는 또 다른 방식으로 한다. 클럭마다 켜진 선 \emph{전부}의 이웃을 한꺼번에 켠다.
세 방식이 같은 답을 내는 근거가 위의 ``순서와 무관하다''이다.

\paragraph{왜 이것이 가설의 검사가 되는가.} 전압을 건 가설이 틀렸고 드럼 위치도 틀리면, 추론은 금세 서로
어긋나 테스트 케이블의 26가닥이 모두 켜진다. 드럼 위치가 맞으면 배선은 모두 참인 추론이 된다. 이때 틀린
가설에서 출발한 전류는 참인 가설의 선에 닿지 못한다. 배선이 양방향이어서, 닿는다면 참인 가설에서 거짓이
유도되는 셈이기 때문이다. 그래서 맞는 위치에서는 테스트 레지스터에 꺼진 선이 남는다.

\section{구동: \texorpdfstring{\code{run\_slice}}{run\_slice}}

\begin{lstlisting}[language=Rust,caption={\code{bombe.rs}: 한 로터 순서의 모든 위치를 돈다(일부).},label={lst:run}]
for g in greeks {
    for l in 0..N as u8 {
        for m in 0..N as u8 {
            let tables = Plugged::tables(order, g, l, m, opts.turnover_at.is_some());
            for r in 0..N as u8 {
                positions += 1;
                plugged.edge_indices(r, opts.turnover_at, &mut idx);
                let live = plugged.propagate(&tables, &idx, start, opts.diagonal_board, true);
                let reg = live[plugged.test_letter as usize];
                if reg != FULL {
                    stops.push(Stop {
                        wheel_order: *order,
                        position: Position { greek: g, left: l, middle: m, right: r },
                        test_register: reg,
                        turnover_at: opts.turnover_at,
                    });
                }
            }
        }
    }
}
\end{lstlisting}

네 겹의 반복문이 드럼 구동 기구다. 가장 안쪽이 오른쪽 로터다. 스크램블러 표(\code{tables})는 가운데 로터의
반복 안, 오른쪽 로터의 반복 \emph{밖}에서 만든다. 앞에서 본 ``바뀌는 것만 다시 계산한다''가 코드의 줄
위치로 나타난 것이다. \code{edge\_indices}는 메뉴의 선마다 26개의 표 가운데 $(r+k) \bmod 26$번째를 고른다.

정지 판정은 한 줄이다. \code{if reg != FULL}. 테스트 레지스터의 26가닥이 전부 켜지지는 않았으면 정지다.
1차 자료가 적은 정지 조건(``입력에 통전되지 않은 점이 있으면 정지'')~\cite[32]{us6812report1944}과 같다.

\paragraph{실제 기계와 다른 점.} 실제 봄베는 에니그마의 왼쪽 로터에 해당하는 드럼을 가장 빨리 돌린다%
~\cite[59]{us6812report1944}. 모델은 오른쪽 로터를 가장 안쪽에 둔다. 모든 위치를 한 번씩 지나기만 하면
정지의 \emph{집합}은 같으므로 표를 다시 쓰기 좋은 순서를 택했다. 회로도 이 순서를 따르게 해서, 정지의
\emph{목록}을 순서까지 비교할 수 있게 했다.

\section{넘어감 가설}

봄베는 크립 구간에서 가운데 로터가 움직이지 않는다고 가정한다. 구간 안에서 가운데 로터가 넘어가면 그 뒤쪽
선들의 스크램블러는 가운데 드럼이 한 칸 틀린 셈이어서, 맞는 위치에서도 정지하지 않는다.

모델은 ``크립 위치 $t$에서 넘어갔다''는 가설을 받을 수 있다(\code{turnover\_at}). 가설이 있으면 표를 두 벌
만든다. 가운데 위치 $m$의 26개와 $m+1$의 26개다. 크립 위치가 $t$ 이상인 선은 두 번째 벌에서 고른다.

\begin{lstlisting}[language=Rust,numbers=none]
out[i] = match turnover_at {
    Some(t) if k >= t => N + r,
    _ => r,
};
\end{lstlisting}

\code{N + r}은 ``두 번째 벌의 \code{r}번째 표''다. 예제 전문을 넘어감이 크립 한가운데에 오는 키로 암호화하면,
맞는 로터 순서를 돌려도 가설 없이는 정지가 하나도 나오지 않고, 가설을 켜면 실제 넘어감 위치의 가설에서
맞는 정지가 나온다.
